\section{Discussion}
\subsection{What the Model Explained}
The model demonstrates how a structured lexical inventory can be coupled to neural retrieval circuits in two processing directions. Its explanatory value lies in making the relation between specification and behavior explicit. Shared graph components influence vector overlap; vector overlap determines match values; the associative-memory variants WTA selection (competition) and thresholded transfer regulate retrieval; and spiking populations approximate the resulting computations over time. The baseline results illustrate this chain for complete, unambiguous input vectors.

A particular strength of the model is that its internal processing can be inspected across all three lexical levels and the associative memories connecting them. The time courses in Section~\ref{sec:activation-interpretation} illustrate how decoded population activity in the conceptual, lemma, and phonological buffers can be examined alongside the selection outputs of the retrieval circuits. This access makes it possible to follow the development of a lexical representation through the processing pathway and to distinguish overlap with related representations from the selection of an additional lexical association. The ability to observe and interpret these internal activity patterns is itself an important outcome of the modeling work: it provides a basis for investigating how representational structure and neural retrieval dynamics jointly produce the observed output.

\subsection{Representational assumptions and emergent behavior}
Similarity is partly built into the model through its relational graphs. It would therefore be circular to treat elevated conceptual or phonological overlap alone as independent confirmation that the model has discovered linguistic similarity. The neural consequences of that structure are a different question. Whether a neighboring match crosses a threshold, survives competition, or alters an output trajectory depends on the retrieval circuit and its dynamics.

The connection between \emph{Tasche} and a vending-machine-related visual node illustrates how individual source associations can influence the resulting similarity structure. Assessing the linguistic adequacy of such associations requires comparison with human semantic judgments. Keeping these associations explicit in the documentation makes their effects traceable and provides a basis for testing alternative lexical specifications. Any revisions can then be recorded and evaluated separately from the behavior of the neural retrieval circuits.

The distinction also applies to the term \emph{coactivation}. Representational overlap, simultaneous selection outputs, and concurrent neuronal firing are different observations. A single distributed representation normally involves many active neurons, and a single selected word can yield several positive vocabulary projections. The geometric baseline in Eq.~\eqref{eq:overlap} makes this distinction testable. It prevents interpretive claims about lexical competition from being based solely on a non-target projection.

\subsection{Relation to earlier work}
The present work builds on the lexical modeling approach investigated in earlier studies of naming, cueing, and production--perception interactions \citep{stille2020,kroeger2016}. The motivation for this methodological account is to make the relationship between lexical data, vector construction, associative-memory operations, and measured neural activity explicit and inspectable. In particular, the paper specifies how within-level similarity arises from the encoded relations, how this differs from cross-level association, and how both affect the interpretation of activation curves. Documentation of the 613-entry inventory and access to the implementation and simulation resources support examination and reuse of these mechanisms. This contribution complements the task-oriented investigations in the earlier publications by providing a formal account of the current lexical core and a basis for subsequent controlled analyses.

The earlier naming study addressed lexical access and cueing within a broader task model \citep{stille2020}. The present account specifies the same basic lexical organization, with conceptual, lemma, and phonological representations linked by associative retrieval. Its narrower scope concerns the tasks used to investigate this organization and the surrounding task architecture; the full 613-entry lexical inventory remains available. The explanatory relationships between representational structure, associative retrieval, and neural activity concern this shared lexical organization. The present study does not claim to reproduce the earlier study's experimental conditions or findings. It also omits the recurrent monitoring architecture considered in earlier production--perception work \citep{kroeger2016}. The current direct concept buffer replaces the additional conceptual-access and cognition stages present in the original WWT implementation developed by Katharina Stille.

The change in random-vector generation, the use of 512 dimensions in the present configuration, and the different surrounding task architecture make the simulations of \citet{stille2020} and the current simulations distinct model realizations. Preserving a recursive construction rule is not equivalent to preserving all vectors, neural weights, or time courses. This distinction permits the present software to serve as a transparent implementation without overstating replication.

\subsection{Comparison with other lexical-access architectures}
\label{sec:lexical-architectures}
The distinction between selecting a lexical identity and retrieving its phonological form connects the present production pathway to established accounts of word production. In WEAVER++, lemma selection precedes word-form encoding, and the selected lemma provides access to its form \citep{levelt1999}. Competition is expressed through a selection rule sensitive to a candidate's activation relative to that of the alternatives. Our production pathway implements a comparable functional division through WTA selection (competition) from concept to lemma, followed by thresholded transfer from lemma to phonological form. The mechanisms differ: WEAVER++ uses a probabilistic selection rule with verification of the relevant lexical links, whereas our first associative memory uses lateral inhibition between neural candidate populations. Moreover, our memories remain continuously connected, so downstream activity follows the evolving lemma state, including selection transients. The timing of these successive transformations therefore emerges from the ongoing neural dynamics.

The interactive two-step model of \citet{dell1997} also separates word selection from phonological retrieval. Its semantic, word, and phoneme units exchange activation through bidirectional excitatory connections. Word selection identifies the most active eligible word, which then receives an activation boost before phoneme selection. This organization allows phonological activation to influence word-level processing during retrieval. Our current implementation places feedback within the WTA selection (competition) circuit, through lateral inhibition, while connections between lexical levels follow the direction of the task. Conceptual and phonological structure enters through precomputed vector geometry, and the output is a complete phonological form vector in production or a concept vector in perception. These choices isolate the contribution of structured representations and associative retrieval; the interactive architecture additionally addresses reciprocal influences among lexical levels and the selection of individual phonemes.

A simpler mathematical reference is the linear heteroassociative map $\mathbf y=VU^{\mathsf T}\mathbf x$, which weights every output vector directly by its input match. Our memories insert the nonlinear selection mechanism $\mathcal F$ between these operations, as shown in Eq.~\eqref{eq:am}. This additional step supplies the thresholding and competition through which the model regulates which lexical associations contribute to the output.

These comparisons distinguish the selection rule, the direction of interlevel coupling, and the unit being retrieved as separate architectural choices. Within the present framework, using the thresholded transfer variant at both transitions would allow every sufficiently strong match to pass at either stage; using the WTA selection (competition) variant at both transitions would add a second competitive selection after lemma retrieval. Adding interlevel feedback would introduce another mechanism, allowing an emerging output representation to influence earlier processing. Such variants could be studied with the same lexicon and input conditions to examine how these choices affect selection dynamics and retrieval under ambiguous input. The current combination assigns a clear role to each memory: WTA selection (competition) of a lemma identity followed by thresholded transfer to its associated target representation.

\subsection{Neural mechanisms and probabilistic interpretation}
\label{sec:neural-probabilistic}
Neural and probabilistic descriptions address complementary aspects of a model. A neural description specifies how neurons and synapses represent information and transform it over time. A probabilistic description specifies probabilities over possible stimuli, interpretations, or responses and the rules for computing or updating them. Both descriptions can apply to the same system: neural population codes can represent probability distributions and support Bayesian inference \citep{ma2006}. Within the NEF, \citet{eliasmithmartens2011} demonstrate probabilistic inference and normalization in a recurrent spiking network. A model can therefore be neural in its realization and probabilistic in the computation it performs.

The present model specifies the dynamics of lexical representations, candidate selection, and associative transfer. Its plotted values have meanings defined by these operations. For example, a thresholded transfer output of 0.8 assigns weight 0.8 to its associated stored vector in Eq.~\eqref{eq:am}. An output near 1 in the WTA selection (competition) memory reflects the configured amplitude of a selected response. Likewise, a positive buffer projection onto \emph{Tasse} during retrieval of \emph{Kanne} can reflect shared vector components, as explained in Section~\ref{sec:activation-interpretation}. Interpreting a value of 0.8 as an 80\% probability that a particular word is intended would introduce an additional assumption about the meaning of that activity.

Such a probabilistic interpretation would start by defining the relevant alternatives and evidence. For lexical recognition, one possible quantity is the probability of each word given an ambiguous input and its context. A model would then specify how lexical expectations and input reliability determine that distribution and how neural activity encodes it. Rescaling activity values to sum to 1 establishes normalization; their interpretation as probabilities additionally depends on these modeling assumptions and their evaluation. The probabilistic selection rule in WEAVER++ discussed in Section~\ref{sec:lexical-architectures} concerns the chance of selecting a candidate \citep{levelt1999}; a probability that a word generated an input concerns lexical uncertainty. These examples illustrate why the meaning of a probability must be specified alongside its neural or algorithmic implementation. The present account focuses on the mechanisms that retrieve and transfer specified lexical representations.

\subsection{Limitations and next analyses}
\label{sec:limitations}
The scope of the present study and the next steps concern several distinct aspects of the model and its evaluation. We address these separately below.

\begin{enumerate}
\item \textbf{Lexical specification and representational scope.} The lexicon is manually specified and includes inherited associations and documented repairs. No lexical learning process is simulated. The lemma layer currently supplies identities without shared grammatical features. Phonological composition does not by itself establish sequence-sensitive encoding. The present tasks isolate lexical retrieval between specified representations; Section~\ref{sec:sensorimotor-interfaces} describes their extension toward acoustic input, articulation, and sound generation.

\item \textbf{Recurrent input maintenance.} A recurrent input buffer is a concrete candidate for the next implementation step (Section~\ref{sec:neural-implementation}). It would add short-term retention to the current input-driven architecture and permit retrieval to be evaluated after a brief cue has disappeared. Implementing this extension would involve specifying how a representation is loaded, held, and cleared or replaced, together with the recurrent gain and synaptic time constant. Evaluation would then examine the stability of the retained vector and its amplitude, as well as retrieval success across different cue durations and retention intervals. Keeping the lexical vectors and associative mappings fixed would allow the contribution of this maintenance mechanism to be assessed against the present sustained-input baseline.

\item \textbf{Coverage of the simulation corpus.} The simulation corpus defined in Section~\ref{sec:simulation-corpus} supports the paper's focus on model description and illustrative retrieval behavior. All 613 lexical entries participate as candidates in the associative memories, and each retrieval decision is evaluated against the full inventory. Input coverage comprises the eight selected words in each direction. This provides sixteen documented examples from the complete implemented lexicon. Extending the corpus with every lexical entry as input in both directions would provide a systematic assessment of retrieval across the inventory.

\item \textbf{Robustness across model configurations.} The reported simulations also use a fixed model configuration: the same parameter settings, vector seed, and network seed are retained across the selected words. This permits a reproducible comparison without item-specific retuning and establishes that one common configuration supports both tasks for the tested inputs. It does not yet quantify sensitivity to alternative random vectors, neural realizations, or parameter choices. The two sources of variation introduced in Sections~\ref{sec:graph-encoding} and~\ref{sec:neural-implementation} can be examined separately: repeating the benchmark across vector seeds would assess dependence on representational geometry, while varying network seeds with the vocabulary fixed would assess variability in its neural implementation. These tests concern robustness beyond the reported realization, rather than the validity of the observed retrieval results themselves.

\item \textbf{Relationship to human behavior.} The relationship to human behavior should likewise be assessed at distinct levels. The separation of concepts, lemmas, and phonological forms, together with the two processing directions, gives the model a psycholinguistically interpretable organization. Earlier work on naming and cueing provides a relevant behavioral context \citep{stille2020}, and the present specification makes the proposed mechanisms available for more detailed examination. Successful retrieval of the selected words nevertheless does not by itself establish quantitative agreement with human reaction times, error distributions, or similarity effects. Such agreement requires explicit comparisons with behavioral observations under corresponding task conditions. The absence of a new behavioral comparison in the present benchmark does not imply that the architecture is unrelated to human processing; it identifies a further step needed to assess the explanatory reach of this particular implementation. Findings from earlier models should inform that assessment while remaining distinguished from validation of the current implementation under the tasks examined here.

\item \textbf{Neuron ablation.} The simulation software also supports neuron ablation in each of the three state buffers and both associative memories of each processing pathway. A selected proportion of neurons can be silenced independently in each component, allowing the effects of reduced neural resources on representation and retrieval to be explored. The \href{https://www.speechtrainer.eu/nef-spa-lexical-model.htm}{online simulation page} provides selected recorded ablation examples; the local simulation supports freely adjustable ablation settings. The present paper documents the intact baseline, while a detailed treatment and systematic evaluation of ablation conditions are reserved for subsequent publications. Earlier insights into such effects are provided by \citet{stille2020}, who examined neuron ablation within lexical buffers and between-level associative memories in a naming task, including the effects of semantic and phonological cues.

\item \textbf{Controlled variation of representational overlap.} A useful next analysis would vary representational overlap while keeping the candidate inventory and neural allocation fixed, and then separate changes in state projections from changes in selection outputs. Repeated vector and network seeds would quantify two different sources of variability. Such experiments would strengthen the distinction between assumed representational structure and consequences of neural retrieval, but they are not needed to identify the mathematical meaning of the quantities documented here.
\end{enumerate}

\subsection{Connecting lexical processing to perception and articulation}
\label{sec:sensorimotor-interfaces}
The tasks in Section~\ref{sec:tasks} examine the lexical component of a broader speech-processing system. In a more complete perception task, an acoustic speech signal would enter an auditory processing pathway, which would construct the phonological representation used for lexical access. In production, the retrieved phonological form provides an intermediate representation on the way to an audible utterance. The broader task requires this form to initiate a temporally organized articulatory program, drive a vocal-tract model, and generate an acoustic output. These extensions would connect the present lexical transformations to observable speech perception and production.

A strength of the NEF-SPA approach is its capacity to integrate these cognitive, sensory, and motor processes. The SPA supplies a common representational format and operations for organizing higher-level processing, while also linking perceptual and motor representations to those computations \citep{eliasmith2013}. The NEF provides methods for realizing representations, transformations, and continuous dynamics in neural populations \citep{eliasmith2003}. Nengo supplies the implementation environment in which these neural populations can be connected to external sensory and motor components, including models of the body and its environment \citep{bekolay2014}. This combination supports the connection of lexical representations to models of auditory input and vocal-tract control within a common computational framework.

Spaun provides an established example of this integration \citep{eliasmith2012}. It receives images of handwritten or printed characters through a simulated eye and generates written responses by controlling a physically modeled arm. Its visual system supplies information for cognitive processing, and its motor system turns the resulting response into an observable movement. The model therefore demonstrates both sides of the connection: visual signals enter through a sensory interface, and answers emerge as writing through a motor effector. For the present speech model, the corresponding objective is to link auditory processing and lexical computation to articulatory movement and acoustic output.

Several components of this route are already realized in our current modeling program. The companion articulatory manuscript documents neural syllable timing, overlapping gesture activation, and the conversion of decoded control signals into vocal-tract movement \citep{kroegerArticulation2026}. An ordered syllable readout for the selected lexical entries provides a shared interface between the lexical and articulatory examples. The current components are evaluated in separate simulations, with symbolic program selection and replay at specified interfaces. A complementary acoustic component is provided by the DYNARTmo extension, which derives vocal-tract acoustics and aerodynamic behavior from articulatory control and geometry and generates speech signals \citep{kroegerAcoustics2026}. Together, these developments provide implemented components spanning lexical retrieval, articulatory control, and sound generation.

The next integration steps are to let accepted lexical outputs initiate the corresponding articulatory programs, connect the neural control stages in a common simulation, and pass their output to the acoustic model. On the perception side, developing an auditory front end that maps an incoming speech signal to the phonological representations required by the lexicon remains a further objective for the present implementation. The recurrent input buffer discussed in Section~\ref{sec:neural-implementation} could then support temporary maintenance of the resulting representation. These extensions would make it possible to evaluate lexical processing together with its sensory input and motor consequences, including the intelligibility of generated speech and the effects of variation in the acoustic input.
